\documentclass[a4paper,9pt]{extarticle}

%=========================================================
% PACKAGES
%=========================================================

\usepackage[
a4paper,
top=1.8cm,
bottom=1.5cm,
footskip=7mm,
left=1cm,
right=1cm,
headheight=14pt,
headsep=5mm
]{geometry}

\usepackage[T1]{fontenc}
\usepackage[utf8]{inputenc}
\usepackage[french]{babel}

\usepackage[sfdefault]{FiraSans}
\renewcommand{\familydefault}{\sfdefault}

\usepackage{multicol}
\setlength{\columnsep}{0.7cm}

\usepackage{amsmath,amssymb}
\usepackage{titlesec}
\usepackage{enumitem}

\usepackage{tcolorbox}
\tcbuselibrary{skins,breakable}

\usepackage{xcolor}
\usepackage{graphicx}
\usepackage{float}

\usepackage{fancyhdr}
\usepackage[hidelinks]{hyperref}


%=========================================================
% PAGE SETTINGS
%=========================================================

\setlength{\parindent}{0pt}
\setlength{\parskip}{1.5pt}

\setlength{\abovedisplayskip}{4pt}
\setlength{\belowdisplayskip}{4pt}

\setlist[itemize]{
leftmargin=*,
itemsep=1pt,
topsep=1pt
}
\setlist[enumerate]{
leftmargin=*,
itemsep=1pt,
topsep=1pt
}


%=========================================================
% COLORS
%=========================================================

\definecolor{Blue}{HTML}{E3F0FF}
\definecolor{Green}{HTML}{E3F7EA}
\definecolor{Pink}{HTML}{FFE5ED}
\definecolor{Orange}{HTML}{FFF0DC}
\definecolor{Purple}{HTML}{EEE7FF}
\definecolor{Yellow}{HTML}{FFFADB}
\definecolor{Goat}{HTML}{F7EEDF}

\definecolor{BlueTitle}{HTML}{285B8F}
\definecolor{GreenTitle}{HTML}{287052}
\definecolor{PinkTitle}{HTML}{A84063}
\definecolor{OrangeTitle}{HTML}{A96500}
\definecolor{PurpleTitle}{HTML}{6245A5}
\definecolor{GoatTitle}{HTML}{8A5A1F}
\definecolor{RecapTitle}{HTML}{5B3A8E}


%=========================================================
% HEADER
%=========================================================

\pagestyle{fancy}
\fancyhf{}

\fancyhead[L]{\small\textbf{Comportement mécanique des matériaux}}
\fancyhead[C]{\small Al--Mg : solution solide, écrouissage, recuit ; Al--Si}
\fancyhead[R]{\small Week 09 --- 02-08/11/2026}

\fancyfoot[C]{\small\href{https://github.com/AlphaXT38G/notes_ba5}{github.com/AlphaXT38G/notes\_ba5}}
\fancyfoot[R]{\small \thepage}

\renewcommand{\headrulewidth}{0.4pt}


%=========================================================
% TITLES
%=========================================================

\titleformat{\section}
{\large\bfseries\color{BlueTitle}}
{}{0pt}{}

\titlespacing{\section}
{0pt}{6pt}{3pt}


%=========================================================
% BOXES
%=========================================================
% Boxes are kept whole (unbreakable): inside multicols, a box split by tcolorbox
% leaves an empty column behind it.

\newtcolorbox{definition}{
enhanced, unbreakable, colback=Blue, colframe=BlueTitle,
boxrule=0.7pt, arc=0mm, title=\textbf{Définition}, fonttitle=\bfseries
}

\newtcolorbox{important}{
enhanced, unbreakable, colback=Pink, colframe=PinkTitle,
boxrule=0.7pt, arc=0mm, title=\textbf{Idée clé}, fonttitle=\bfseries
}

\newtcolorbox{formula}{
enhanced, unbreakable, colback=Green, colframe=GreenTitle,
boxrule=0.7pt, arc=0mm, title=\textbf{Formules clés}, fonttitle=\bfseries
}

\newtcolorbox{warning}{
enhanced, unbreakable, colback=Orange, colframe=OrangeTitle,
boxrule=0.7pt, arc=0mm, title=\textbf{Attention}, fonttitle=\bfseries
}

\newtcolorbox{example}{
enhanced, unbreakable, colback=Purple, colframe=PurpleTitle,
boxrule=0.7pt, arc=0mm, title=\textbf{Exemple}, fonttitle=\bfseries
}

% GOAT - brown: the idea told as simply as possible
\newtcolorbox{goat}{
enhanced, unbreakable, colback=Goat, colframe=GoatTitle,
boxrule=0.7pt, arc=0mm, title=\textbf{Expliqué à une chèvre}, fonttitle=\bfseries
}

% RECAP - violet: what to know for the exam, at the end of the week
\newtcolorbox{recap}{
enhanced, breakable, colback=Purple, colframe=RecapTitle,
boxrule=1pt, arc=0mm, title=\textbf{Récap pour l'examen}, fonttitle=\bfseries
}

% Screenshot of a slide of the course: \slide{cNN-PP}
\newcommand{\slide}[2][0.95]{%
\begin{center}
\includegraphics[width=#1\linewidth]{../images/course/#2}
\end{center}}


%=========================================================
% DOCUMENT
%=========================================================

\begin{document}

\thispagestyle{fancy}

\begin{multicols}{2}

%=========================================================
\section{Al--Mg : solution solide et écrouissage}

\begin{important}
Dans Al--Mg les précipités durcissent peu. On utilise :
\begin{itemize}
\item \textbf{solution solide} : champ de contraintes du soluté (gros : compression ; petit :
      traction) et de la dislocation interagissent ; attraction \emph{ou} répulsion, les deux
      freinent ;
\item \textbf{écrouissage} : les dislocations se gênent entre elles, et la déformation les
      multiplie.
\end{itemize}
\end{important}

\slide[0.75]{c03b-58}

\begin{goat}
Tout ce qui gêne les dislocations durcit le métal : des atomes d'une autre taille, ou les autres
dislocations qu'on crée en le déformant.
\end{goat}

%=========================================================
\section{Recuit : restauration, recristallisation}

\begin{important}
5--10\,\% du travail de déformation reste stocké dans les dislocations ; chauffer les élimine
(\textbf{recuit}) :
\begin{itemize}
\item \textbf{restauration} : annihilation, réarrangement ; adoucissement modeste ;
\item \textbf{recristallisation} : nouveaux grains \textbf{équiaxes} presque sans dislocations ;
      seuil de déformation ; plus d'écrouissage $\Rightarrow$ grains plus fins.
\end{itemize}
\end{important}

\slide[0.85]{c03b-66}

\begin{important}
Joints mobiles à chaud : \textbf{croissance de grain} (réduire l'aire, comme une mousse),
parfois \textbf{anormale}. Fins précipités \textbf{épinglent} les joints. Déformation + recuit
partiel > écrouissage seul.
\end{important}

\begin{goat}
Chauffer un métal écroui permet de démêler les dislocations ou de refaire des grains neufs : il
redevient mou.
\end{goat}

%=========================================================
\section{Al--Si et désignations}

\begin{important}
Al--Si : alliages de \textbf{fonderie}. Retrait $\Rightarrow$ retassures, \textbf{masselottes} ;
H dissous $\Rightarrow$ bulles, \textbf{dégazage}. Si en plaquettes ; \textbf{modification}
(Sr, Na ou solidification rapide) : Si fin et ramifié, meilleures propriétés.
\end{important}

\begin{definition}
Corroyés 4 chiffres, fonderie 3. 1xxx pur, 2xxx Cu, 3xxx Mn, 4xxx Si, 5xxx Mg, 6xxx Mg+Si,
7xxx Zn, 8xxx autres. \textbf{Durcissables} : 2xxx (Al$_2$Cu), 6xxx (Mg$_2$Si), 7xxx (MgZn$_2$).
\end{definition}

\begin{goat}
Pour couler, l'aluminium au silicium ; les familles 2, 6, 7 se durcissent par traitement
thermique.
\end{goat}

%=========================================================
\section{Récap de la semaine}

\begin{recap}
\begin{itemize}
\item Trois mécanismes (avec la semaine 08) : écrouissage (obstacle : dislocations), solution
      solide (atomes étrangers), durcissement structural (précipités).
\item Recuit : restauration, recristallisation (seuil, grain fin si très déformé), croissance
      normale/anormale, épinglage.
\item Al--Si : retrait, masselottes, dégazage, modification.
\item Séries ; durcissables 2xxx, 6xxx, 7xxx.
\end{itemize}
\end{recap}

\end{multicols}

\end{document}
